\documentclass[11pt]{article}

\usepackage[utf8]{inputenc}
\usepackage[T1]{fontenc}
\usepackage[spanish,es-tabla]{babel}
\usepackage[a4paper,margin=2.2cm]{geometry}
\usepackage{booktabs}
\usepackage{tabularx}
\usepackage{float}
\usepackage{xcolor}
\usepackage{hyperref}
\usepackage{amsmath}

\hypersetup{
  colorlinks=true,
  linkcolor=blue!55!black,
  urlcolor=blue!55!black
}

\newcommand{\img}[2]{%
  \begingroup
  \pdfximage width #1 {#2}%
  \pdfrefximage\pdflastximage
  \endgroup
}

\title{Resumen comparativo: Chua no suave y Chua arctan}
\author{Pipeline Nyquist/DF--EFORK--verificacion de ocultamiento}
\date{8 de mayo de 2026}

\begin{document}
\maketitle

\section*{Estado de las corridas}

Se interrumpio la continuidad de trabajo sobre la etapa de bifurcacion y se
consolido el resumen de los dos ejemplos disponibles. Para el caso no suave se
dispone de una corrida operativa con verificacion de ocultamiento, cuencas y
figuras de trayectorias desde vecindades de equilibrio. Para el caso
\texttt{arctan}, la corrida con el mismo contrato numerico produjo las figuras
iniciales, pero la continuacion divergio antes de poder construir una referencia
robusta para el atractor candidato.

\begin{table}[H]
\centering
\caption{Comparacion operativa de las corridas en \(q=0.99\).}
\begin{tabularx}{\textwidth}{lXX}
\toprule
Concepto & Chua no suave & Chua \texttt{arctan} \\
\midrule
Modelo & \texttt{piecewise} & \texttt{arctan} \\
Carpeta de salida &
\begin{tabular}[t]{@{}l@{}}\texttt{chua\_no\_suave\_q0p99\_run}\\\texttt{chua\_piecewise}\end{tabular} &
\begin{tabular}[t]{@{}l@{}}\texttt{chua\_arctan\_run}\\\texttt{chua\_arctan}\end{tabular} \\
Orden fraccionario & \(q=0.99\) & \(q=0.99\) \\
Contrato EFORK & \(h=0.01,\ L_m=5,\ T_{\rm tr}=80,\ T_{\rm keep}=100\) &
\(h=0.01,\ L_m=5,\ T_{\rm tr}=80,\ T_{\rm keep}=100\) \\
Semilla final de continuacion &
\begin{tabular}[t]{@{}l@{}}\((2.088159,\ 1.540769,\)\\\(-1.505487)\)\end{tabular} &
\begin{tabular}[t]{@{}l@{}}\((-1.3415{\times}10^{10},\)\\\(-2.1666{\times}10^{10},\)\\\(-2.3332{\times}10^{10})\)\end{tabular} \\
Verificacion oculta &
Ejecutada; 2 trayectorias \texttt{TARGET} de 72 muestras. &
No valida; el backend no pudo construir referencia porque la semilla final divergio. \\
Lectura principal &
Hay candidato acotado, pero no es oculto fuerte: existe conexion desde vecindad de \(E_+\). &
No hay atractor candidato valido con este contrato; primero hay que estabilizar la continuacion. \\
\bottomrule
\end{tabularx}
\end{table}

\section*{Chua no suave}

La corrida no suave genero el atractor candidato, la seccion de referencia, el
resumen de sondas desde equilibrios, la ilustracion de ocultamiento y las
cuencas. En la verificacion oculta se usaron tres equilibrios y tres radios:
\(10^{-4}\), \(10^{-3}\) y \(10^{-2}\), con 8 muestras por radio.

\begin{table}[H]
\centering
\caption{Conteos de la verificacion de ocultamiento para Chua no suave.}
\begin{tabular}{llrrrrr}
\toprule
Equilibrio & Radio & EQ & DIV & TARGET & OTHER & UNKNOWN \\
\midrule
\(E_0\) & \(10^{-4}\) & 8 & 0 & 0 & 0 & 0 \\
\(E_0\) & \(10^{-3}\) & 8 & 0 & 0 & 0 & 0 \\
\(E_0\) & \(10^{-2}\) & 8 & 0 & 0 & 0 & 0 \\
\(E_+\) & \(10^{-4}\) & 8 & 0 & 0 & 0 & 0 \\
\(E_+\) & \(10^{-3}\) & 8 & 0 & 0 & 0 & 0 \\
\(E_+\) & \(10^{-2}\) & 0 & 6 & 2 & 0 & 0 \\
\(E_-\) & \(10^{-4}\) & 8 & 0 & 0 & 0 & 0 \\
\(E_-\) & \(10^{-3}\) & 8 & 0 & 0 & 0 & 0 \\
\(E_-\) & \(10^{-2}\) & 0 & 5 & 0 & 3 & 0 \\
\bottomrule
\end{tabular}
\end{table}

El resultado clave es que aparecen 2 trayectorias \texttt{TARGET}, ambas desde
la vecindad de \(E_+\) con radio \(10^{-2}\). Por tanto, el caso no suave no
queda certificado como atractor oculto fuerte bajo este muestreo: el atractor
candidato tiene conexion detectable desde una vecindad de equilibrio.

\begin{figure}[H]
\centering
\img{0.72\textwidth}{../chua_no_suave_q0p99_run/chua_piecewise/final_pdf_figs/fig03_final_attractor.png}
\caption{Atractor candidato para el Chua no suave en \(q=0.99\).}
\end{figure}

\begin{figure}[H]
\centering
\img{0.82\textwidth}{../chua_no_suave_q0p99_run/chua_piecewise/final_pdf_figs/fig05b_hiddenness_overview.png}
\caption{Atractor candidato, equilibrios y trayectorias clasificadas desde vecindades de equilibrio.}
\end{figure}

\begin{figure}[H]
\centering
\img{0.72\textwidth}{../chua_no_suave_q0p99_run/chua_piecewise/final_pdf_figs/fig06_basin_overlay.png}
\caption{Cuenca operacional en el plano usado por el pipeline para el Chua no suave.}
\end{figure}

\section*{Chua arctan}

La corrida \texttt{arctan} se ejecuto con el mismo orden fraccionario y el mismo
perfil rapido, sin bifurcaciones. Nyquist/DF encontro dos cruces y selecciono la
rama 0:
\[
  \omega_0 = 2.0991692817,\qquad k=-1.3282768771,\qquad a_0=0.9019225739.
\]
Sin embargo, al continuar hasta \(\epsilon=1\), el estado final fue de orden
\(10^{10}\):
\[
(-1.3415{\times}10^{10},\ -2.1666{\times}10^{10},\ -2.3332{\times}10^{10}).
\]
Por ello, el backend de ocultamiento no pudo construir la seccion de referencia
del atractor objetivo. La conclusion para este contrato es que el caso
\texttt{arctan} no proporciona un atractor candidato valido; requiere cambiar
contrato numerico, rama, parametros de \(\arctan\), o semilla.

\begin{figure}[H]
\centering
\img{0.72\textwidth}{../chua_arctan_run/chua_arctan/final_pdf_figs/fig01_nyquist_df.png}
\caption{Cruces Nyquist/DF para la corrida \texttt{arctan}.}
\end{figure}

\begin{figure}[H]
\centering
\img{0.72\textwidth}{../chua_arctan_run/chua_arctan/final_pdf_figs/fig02_continuation_progress.png}
\caption{Progreso de continuacion para \texttt{arctan}; la rama seleccionada termina divergiendo.}
\end{figure}

\begin{figure}[H]
\centering
\img{0.72\textwidth}{../chua_arctan_run/chua_arctan/final_pdf_figs/fig03_final_attractor.png}
\caption{Figura final generada para \texttt{arctan}. No debe interpretarse como atractor oculto validado, porque la verificacion fallo por divergencia.}
\end{figure}

\section*{Conclusion}

Entre las dos corridas, el caso no suave es el unico que produjo un candidato
acotado con verificacion posterior. Aun asi, no es un candidato oculto fuerte:
la muestra detecto trayectorias que salen de \(E_+\) y alcanzan el atractor
objetivo. El caso \texttt{arctan}, con el contrato usado aqui, no llega a la fase
de verificacion porque la continuacion diverge.

El siguiente paso razonable es no continuar con bifurcaciones sino probar una
de estas alternativas: cambiar a la rama Nyquist/DF 1, reducir el paso/ajustar
la memoria para \texttt{arctan}, o hacer un barrido de parametros \((a_1,a_2,\rho)\)
antes de repetir la verificacion oculta.

\end{document}
